% ==========================================
% SECCIÓN PRELIMINAR: RESUMEN
% ==========================================

\pagestyle{empty}
\onehalfspacing
\fontsize{12}{18}\selectfont

\section*{RESUMEN}

La gestión académico-administrativa de los docentes becarios de doctorado en la Universidad Central del Ecuador (UCE) enfrenta una problemática crítica derivada del uso de hojas de cálculo aisladas y almacenamiento local. Esta práctica ocasiona fragmentación de datos, falta de alertas automáticas de vencimiento, dificultades para emitir certificados y un alto riesgo de error humano en el control de expedientes, obligando a los docentes a realizar visitas presenciales recurrentes solo para consultar qué documentos tienen entregados o pendientes.

Para resolver esta vulnerabilidad, se propone el diseño e implementación de una plataforma web modular desacoplada denominada Sistema de Gestión de Doctorados (SGD), interconectada mediante una API REST en Express.js con base de datos PostgreSQL y Prisma ORM. La solución integra un Panel Administrativo en Next.js 16 (React 19) con control RBAC para que el personal de la Unidad gestione integralmente programas doctorales, expedientes y la agenda de citas sin depender de Excel. Paralelamente, incorpora un Portal Estudiantil en Astro y React donde los docentes becarios consultan con transparencia el estado de su expediente (documentos entregados y faltantes), agendan citas en línea en tiempo real con los miembros de la Unidad para atención personalizada, y autogestionan solicitudes, alertas y certificados digitales. La plataforma se evalúa mediante una demo funcional en la nube sobre un banco de pruebas sintético con perfiles ficticios estructurado según formatos oficiales bajo la LOPDP, respaldada por manuales interactivos en Fumadocs.

\vspace{0.5cm}
\textbf{Palabras clave:} Sistema de Gestión de Doctorados (SGD), Arquitectura Desacoplada, Next.js 16, Astro, Prisma ORM, Control RBAC, Reservas de Citas en Tiempo Real, Fumadocs, UCE.

\newpage
